\documentclass[12pt,a4paper]{article}
\usepackage[UTF8]{ctex}
\usepackage{amsmath,amssymb,amsthm}
\usepackage{geometry}

\geometry{margin=2.5cm}

\title{稳态解与稳定性：标量与矩阵情况的对应}
\author{第11章补充材料}
\date{}

\newtheorem{theorem}{定理}
\newtheorem{definition}{定义}
\newtheorem{lemma}{引理}
\newtheorem{proof}{证明}
\newtheorem{corollary}{推论}

\begin{document}

\maketitle

\section{问题陈述}

考虑两种微分方程：

\textbf{标量情况：}
\begin{equation}
\dot{x}(t) = ax(t), \quad x(0) = x_0,
\end{equation}
其中$a \in \mathbb{R}$是常数。

\textbf{向量情况：}
\begin{equation}
\dot{\mathbf{x}}(t) = A\mathbf{x}(t), \quad \mathbf{x}(0) = \mathbf{x}_0,
\end{equation}
其中$\mathbf{x}(t) \in \mathbb{R}^n$，$A \in \mathbb{R}^{n \times n}$是常数矩阵。

我们要证明和理解：
\begin{enumerate}
\item 标量情况：如果$a < 0$，则$x = 0$是稳态解（平衡点），且从任何初始条件$x_0$，解都收敛到$x = 0$。
\item 向量情况：如果矩阵$A$的所有特征值具有负实部（$A$是"负定"的，更准确地说，$A$是Hurwitz稳定的），则$\mathbf{x} = \mathbf{0}$是稳态解，且从任何初始条件$\mathbf{x}_0$，解都收敛到$\mathbf{x} = \mathbf{0}$。
\end{enumerate}

\section{平衡点（稳态解）的定义}

\begin{definition}[平衡点]
对于微分方程$\dot{x} = f(x)$，如果$f(x^*) = 0$，则$x^*$称为平衡点（或稳态解）。
\end{definition}

平衡点的物理意义是：如果系统从平衡点开始，它将永远停留在该点，因为在该点处变化率为零。

\section{标量情况的完整分析}

\subsection{平衡点的存在性}

对于标量方程(1)，平衡点满足：
\begin{equation}
\dot{x} = ax = 0.
\end{equation}

如果$a \neq 0$，则唯一的平衡点是$x^* = 0$。

如果$a = 0$，则所有点都是平衡点（系统是静态的）。

\subsection{解的显式形式}

标量方程(1)的解为：
\begin{equation}
x(t) = e^{at}x_0.
\end{equation}

\subsection{稳定性分析}

\begin{theorem}[标量情况的稳定性]
对于标量微分方程$\dot{x} = ax$，平衡点$x^* = 0$的稳定性如下：
\begin{enumerate}
\item 如果$a < 0$，则平衡点$x^* = 0$是\textbf{渐近稳定}的：从任何初始条件$x_0$，解$x(t) = e^{at}x_0 \to 0$当$t \to \infty$。
\item 如果$a > 0$，则平衡点$x^* = 0$是\textbf{不稳定}的：从任何非零初始条件$x_0 \neq 0$，解$|x(t)| = |e^{at}x_0| \to \infty$当$t \to \infty$。
\item 如果$a = 0$，则平衡点是\textbf{临界稳定}的：解$x(t) = x_0$保持常数。
\end{enumerate}
\end{theorem}

\begin{proof}
根据解的形式$x(t) = e^{at}x_0$：

\textbf{情况1：$a < 0$}
\begin{align}
\lim_{t \to \infty} x(t) &= \lim_{t \to \infty} e^{at}x_0 \\
&= x_0 \lim_{t \to \infty} e^{at} \\
&= x_0 \cdot 0 = 0.
\end{align}
因此，无论初始条件$x_0$如何，解都收敛到平衡点$x^* = 0$。这证明了平衡点是渐近稳定的。

\textbf{情况2：$a > 0$}
\begin{align}
\lim_{t \to \infty} |x(t)| &= \lim_{t \to \infty} |e^{at}x_0| \\
&= |x_0| \lim_{t \to \infty} e^{at} \\
&= |x_0| \cdot \infty = \infty \quad \text{（如果$x_0 \neq 0$）}.
\end{align}
因此，从任何非零初始条件，解都发散，平衡点是不稳定的。

\textbf{情况3：$a = 0$}
\begin{equation}
x(t) = e^{0 \cdot t}x_0 = x_0,
\end{equation}
解保持常数，既不收敛也不发散。
\end{proof}

\subsection{直观理解}

当$a < 0$时：
\begin{itemize}
\item 如果$x > 0$，则$\dot{x} = ax < 0$，系统会减小$x$的值。
\item 如果$x < 0$，则$\dot{x} = ax > 0$（因为$a < 0$且$x < 0$），系统会增加$x$的值。
\item 因此，系统总是被"拉向"平衡点$x = 0$。
\end{itemize}

这就像一个有阻尼的系统，总是倾向于回到平衡位置。

\section{向量情况的完整分析}

\subsection{平衡点的存在性}

对于向量方程(2)，平衡点满足：
\begin{equation}
\dot{\mathbf{x}} = A\mathbf{x} = \mathbf{0}.
\end{equation}

如果矩阵$A$是可逆的，则唯一的平衡点是$\mathbf{x}^* = \mathbf{0}$。

如果$A$是奇异的（不可逆），则存在非零的平衡点（零空间中的点）。

\subsection{解的显式形式}

向量方程(2)的解为：
\begin{equation}
\mathbf{x}(t) = e^{At}\mathbf{x}_0.
\end{equation}

\subsection{特征值与稳定性}

矩阵$A$的稳定性由其特征值决定。设$A$的特征值为$\lambda_1, \lambda_2, \ldots, \lambda_n$（可能是复数）。

\begin{definition}[Hurwitz稳定矩阵]
如果矩阵$A$的所有特征值都具有负实部，即$\text{Re}(\lambda_i) < 0$对所有$i$成立，则称$A$是Hurwitz稳定的。
\end{definition}

\textbf{注意：}这里"负定"的说法不够准确。严格来说：
\begin{itemize}
\item \textbf{负定矩阵}：对于对称矩阵$A$，如果$\mathbf{x}^T A \mathbf{x} < 0$对所有非零$\mathbf{x}$成立，则$A$是负定的。
\item \textbf{Hurwitz稳定矩阵}：如果$A$的所有特征值具有负实部，则$A$是Hurwitz稳定的。
\item 对于对称矩阵，负定$\Leftrightarrow$所有特征值为负$\Leftrightarrow$Hurwitz稳定。
\item 对于非对称矩阵，Hurwitz稳定是更一般的概念。
\end{itemize}

\subsection{稳定性定理}

\begin{theorem}[向量情况的稳定性]
对于向量微分方程$\dot{\mathbf{x}} = A\mathbf{x}$，平衡点$\mathbf{x}^* = \mathbf{0}$的稳定性如下：
\begin{enumerate}
\item 如果矩阵$A$的所有特征值具有负实部（$A$是Hurwitz稳定的），则平衡点$\mathbf{x}^* = \mathbf{0}$是\textbf{渐近稳定}的：从任何初始条件$\mathbf{x}_0$，解$\mathbf{x}(t) = e^{At}\mathbf{x}_0 \to \mathbf{0}$当$t \to \infty$。
\item 如果矩阵$A$至少有一个特征值具有正实部，则平衡点$\mathbf{x}^* = \mathbf{0}$是\textbf{不稳定}的。
\item 如果矩阵$A$的所有特征值具有非正实部，且具有零实部的特征值对应的Jordan块都是一阶的，则平衡点是\textbf{临界稳定}的。
\end{enumerate}
\end{theorem}

\begin{proof}
假设$A$可以对角化（对于不能对角化的情况，可以使用Jordan标准型，证明类似）：
\begin{equation}
A = P\Lambda P^{-1},
\end{equation}
其中$\Lambda = \text{diag}(\lambda_1, \lambda_2, \ldots, \lambda_n)$。

则解为：
\begin{equation}
\mathbf{x}(t) = e^{At}\mathbf{x}_0 = Pe^{\Lambda t}P^{-1}\mathbf{x}_0 = P\text{diag}(e^{\lambda_1 t}, e^{\lambda_2 t}, \ldots, e^{\lambda_n t})P^{-1}\mathbf{x}_0.
\end{equation}

\textbf{情况1：所有特征值具有负实部}

对于每个特征值$\lambda_i$，有$\text{Re}(\lambda_i) < 0$，因此：
\begin{equation}
|e^{\lambda_i t}| = e^{\text{Re}(\lambda_i) t} \to 0 \quad \text{当} \quad t \to \infty.
\end{equation}

因此：
\begin{equation}
\text{diag}(e^{\lambda_1 t}, e^{\lambda_2 t}, \ldots, e^{\lambda_n t}) \to \mathbf{0} \quad \text{（零矩阵）},
\end{equation}
从而：
\begin{equation}
\mathbf{x}(t) = P\text{diag}(e^{\lambda_1 t}, e^{\lambda_2 t}, \ldots, e^{\lambda_n t})P^{-1}\mathbf{x}_0 \to P \cdot \mathbf{0} \cdot P^{-1}\mathbf{x}_0 = \mathbf{0}.
\end{equation}

这证明了平衡点$\mathbf{x}^* = \mathbf{0}$是渐近稳定的。

\textbf{情况2：至少有一个特征值具有正实部}

设$\lambda_j$具有正实部，即$\text{Re}(\lambda_j) > 0$。选择初始条件$\mathbf{x}_0 = P\mathbf{e}_j$（其中$\mathbf{e}_j$是第$j$个标准基向量），则：
\begin{equation}
\mathbf{x}(t) = Pe^{\Lambda t}P^{-1}P\mathbf{e}_j = Pe^{\Lambda t}\mathbf{e}_j = P(e^{\lambda_j t}\mathbf{e}_j) = e^{\lambda_j t}P\mathbf{e}_j.
\end{equation}

由于$\text{Re}(\lambda_j) > 0$，有$|e^{\lambda_j t}| = e^{\text{Re}(\lambda_j) t} \to \infty$，因此$\|\mathbf{x}(t)\| \to \infty$，平衡点是不稳定的。
\end{proof}

\section{标量与向量情况的对应关系}

\subsection{对应表}

\begin{table}[h]
\centering
\begin{tabular}{|c|c|}
\hline
\textbf{标量情况} & \textbf{向量情况} \\
\hline
$a \in \mathbb{R}$ & $A \in \mathbb{R}^{n \times n}$ \\
\hline
特征值：$a$本身 & 特征值：$\lambda_1, \lambda_2, \ldots, \lambda_n$ \\
\hline
$a < 0$ & 所有$\text{Re}(\lambda_i) < 0$ \\
\hline
$a > 0$ & 至少一个$\text{Re}(\lambda_i) > 0$ \\
\hline
$a = 0$ & 所有$\text{Re}(\lambda_i) \leq 0$，且零实部特征值对应一阶Jordan块 \\
\hline
解：$x(t) = e^{at}x_0$ & 解：$\mathbf{x}(t) = e^{At}\mathbf{x}_0$ \\
\hline
平衡点：$x^* = 0$ & 平衡点：$\mathbf{x}^* = \mathbf{0}$ \\
\hline
稳定性：由$a$的符号决定 & 稳定性：由$A$的特征值实部决定 \\
\hline
\end{tabular}
\caption{标量与向量情况的对应关系}
\end{table}

\subsection{关键对应}

\textbf{标量情况：}
\begin{itemize}
\item 参数$a$既是系数，也是唯一的"特征值"
\item $a < 0$意味着系统有负的"增长率"，因此解衰减到零
\end{itemize}

\textbf{向量情况：}
\begin{itemize}
\item 矩阵$A$有多个特征值$\lambda_i$
\item 所有$\text{Re}(\lambda_i) < 0$意味着所有"增长模式"都是衰减的，因此解衰减到零
\item 即使$A$不是对称的，"负定"的概念也通过特征值的负实部来体现
\end{itemize}

\section{直观理解}

\subsection{标量情况：$a < 0$}

想象一个弹簧-质量系统，有很强的阻尼：
\begin{itemize}
\item 如果质量偏离平衡位置（$x \neq 0$），阻尼力会将其拉回
\item 阻尼系数越大（$|a|$越大），回到平衡位置越快
\item 最终，系统会稳定在平衡位置$x = 0$
\end{itemize}

\subsection{向量情况：$A$的所有特征值具有负实部}

想象一个多自由度系统：
\begin{itemize}
\item 每个特征值对应系统的一个"模式"
\item 如果所有模式的"增长率"都是负的（所有特征值具有负实部），则所有模式都会衰减
\item 无论初始条件如何，系统最终都会稳定在平衡点$\mathbf{x} = \mathbf{0}$
\end{itemize}

\section{几何解释}

\subsection{标量情况}

在相空间（一维）中：
\begin{itemize}
\item 平衡点$x = 0$是一个点
\item 如果$a < 0$，所有轨迹都指向这个点（吸引子）
\item 如果$a > 0$，所有轨迹都远离这个点（排斥子）
\end{itemize}

\subsection{向量情况}

在相空间（$n$维）中：
\begin{itemize}
\item 平衡点$\mathbf{x} = \mathbf{0}$是一个点
\item 如果$A$的所有特征值具有负实部，所有轨迹都指向这个点（稳定节点或稳定焦点）
\item 如果$A$有正实部特征值，至少有一个方向上的轨迹远离这个点（不稳定）
\end{itemize}

\section{实际应用：误差动力学}

在机器人控制中，我们关心误差动力学：
\begin{equation}
\dot{\theta}_e = -K_p \theta_e,
\end{equation}
其中$K_p > 0$是控制增益。

这对应标量情况，其中$a = -K_p < 0$，因此误差$\theta_e$会收敛到零。

对于多关节系统：
\begin{equation}
\dot{\mathbf{x}}_e = A\mathbf{x}_e,
\end{equation}
如果设计控制器使得$A$的所有特征值具有负实部，则误差$\mathbf{x}_e$会收敛到零。

\section{结论}

我们证明了：

\begin{enumerate}
\item \textbf{标量情况}：如果$a < 0$，则$x = 0$是渐近稳定的平衡点，从任何初始条件，解都收敛到$x = 0$。

\item \textbf{向量情况}：如果矩阵$A$的所有特征值具有负实部（$A$是Hurwitz稳定的），则$\mathbf{x} = \mathbf{0}$是渐近稳定的平衡点，从任何初始条件，解都收敛到$\mathbf{x} = \mathbf{0}$。

\item \textbf{对应关系}：标量情况是向量情况的特例（$n = 1$）。标量的符号对应矩阵特征值的实部符号。

\item \textbf{物理意义}：负的"增长率"（标量）或负的特征值实部（矩阵）意味着系统具有"恢复力"，总是倾向于回到平衡位置。
\end{enumerate}

这些结果构成了线性系统稳定性理论的基础，在控制系统的设计和分析中至关重要。

\end{document}

